\scanpage{25}
\mspage{24}

En effet, $\widetilde{F\vee G}=\widetilde F\cup\widetilde G$, et $\widetilde F$ et $\widetilde G$ y sont fermés, \deleted{\unsure{dans|de} $\unsure{M|\mathcal{M}}$} ce qui prouve que
\[
 \mathrm{Multomb}(F\cup G)=\mathrm{Multomb}(F)\cup\mathrm{Mulomb}(\unsure{C|G}),
\]
et que $\mathrm{Mulomb}(F)$ et $\mathrm{Mulomb}(G)$ y sont fermés, donc \unsure{(*)} et la première relation \unsure{(**)|($\ll$)}. La deuxième relation \inserted{\unsure{(**)|($\ll$)}} provient de $\unsure{=|\eqsim|\overline{m}}$ de $\widetilde{F\cap G}=\widetilde F\cap\widetilde G$.

\uline{NB} Il n'est pas vrai en général que $\mathrm{Mulomb}(F)\between\mathrm{Mulomb}(G)$ implique $F\between\unsure{C|G}$, même si $F$ et $G$ sont \unsure{réduits|s'écrivent} aux lieux distincts $x$ et $y$ --- auquel cas $\mathrm{Multomb}(X)=\{\{\unsure{\lambda|x}\}\}$, $\mathrm{Mulomb}(Y)=\{\{y\}\}$ sont compatibles et $\unsure{=|\eqsim}$ \unsure{disjointes|disjoints}, mais $x$ et $y$ ne le sont pas forcément. Mais il \unsure{ser.|sera} vrai \unsure{par contre qu'on a équivalence|pour d'autres quasi-équivalences|pour d'autre quasi-équivalence}, dans le cas où $\unsure{\mathcal{F}|\mathfrak{F}|\mathscr{F}}$ est \unsure{un|une} atelier \unsure{ensembliste|ens.embliste}\unsure{.|:}
\margin{\unsure{vérifier|vérifies} plus tard}

\uline{Corollaire 2} Si $F$ et $G$ sont \unsure{disjoints|disjointes|disjoint}, alors $\unsure{\mathrm{Mul.mb}|\mathrm{Mulomb}}(F)$ et $\mathrm{Mulomb}(G)$ aussi.

Cela provient formellement du corollaire 1 et de la formule
\[
 \mathrm{Mulomb}(\unsure{\varnothing|\phi}_{\unsure{\mathcal{F}|F}})=\varnothing\quad(\text{figure vide dans }\unsure{M|\mathcal M})
 \tag*{(1.40)}
\]

\uline{Proposition} Soient $\unsure{F|\mathbb{F}},G\in\mathcal{F}$. Alors $F\ll G$
\[
 \Longleftrightarrow\quad\mathrm{Mulomb}(F)\ll\mathrm{Mulomb}(G).
\]
